

\lussec 6. On-shell O($\mib a$) improvement 

The lattice-spacing dependence 
of correlation functions involving the fields at non-zero flow time
is best discussed in the theory in 4+1 dimensions.
In this framework, the standard argumentation 
based on locality, power counting and symmetry
can be applied. 
The relevant Symanzik local effective theory
[\cite{SymanzikSeffI},\cite{SymanzikSeffII}] then
coincides with the continuum theory in 4+1 dimensions,
with O($a$) and higher-order terms in the lattice spacing $a$
added to the action and the local fields.
Moreover, on-shell improvement [\cite{OnShell}]
means that the cancellation of lattice-spacing effects
is limited to correlation functions 
at non-zero distances in 4+1 dimensions.

Although the theoretical discussion is more widely applicable,
the O($a$)-improved Wilson formulation of lattice QCD
[\cite{SW},\cite{SFimp}] is assumed from now on. 
Where possible, the same conventions
and notation as in ref.~[\cite{SFimp}] are used.


\lussubsec 6.1 Effective action

Since the flow equations in the lattice theory 
are classically O($a$) improved,
and since loop diagrams do not contribute to the correlation 
functions at asymptotically large flow times,
all O($a$) terms in the Symanzik effective action 
must be boundary terms, i.e.~local terms at flow time zero.
Furthermore, if the effective theory is to describe on-shell
quantities only,
as is the case here, many terms can be eliminated using the 
quark field equations,
the flow equations and the boundary conditions (2.1),(2.6). 
In the chiral limit, a possible choice 
of the O($a$) term in the effective action is then
\equation{
  a\int\rmd^4x\left\{c_1\psibar(x)\frac{i}{4}\sigma_{\mu\nu}F_{\mu\nu}(x)\psi(x)
  +c_2\lambdabar(0,x)\lambda(0,x)\right\},
  \enum
} 
where $F_{\mu\nu}(x)$ is the field tensor of the fundamental
gauge potential 
(here and below, $c_1,c_2,\ldots$ denote coefficients of the 
effective theory).

Quite many more terms are required to describe the 
O($a$) lattice-spacing effects at non-zero quark masses.
In the theory in 4 dimensions,
these were completely classified
by Bhattacharya et al.~[\cite{ImpNonD}].
The Symanzik effective action 
in 4+1 dimensions in addition includes a term
\equation{
  a\int\rmd^4x\left\{
  \lambdabar(0,x)(c_3M+c_4\kern0.5pt\tr M)\psi(x)+
  \psibar(x)(c_3M+c_4\kern0.5pt\tr M)\lambda(0,x)\right\}
  \enum
}
that contributes to correlation functions involving
the bulk fields, $M$ being the quark mass matrix in the effective theory. 
Its effect on the correlation functions is equivalent
to a renormalization
\equation{
  \chi(t,x)\to Z\chi(t,x),
  \quad
  \lambda(t,x)\to Z^{-1}\lambda(t,x),
  \enum
  \nexteq{2.5ex}
  Z=1+ac_3M+ac_4\kern0.5pt\tr M,
  \enum
}
of the bulk fermion (and anti-fermion) fields. 


\lussubsec 6.2 Effective local fields

The form of the effective fields representing
the local lattice fields in the Symanzik theory
is constrained by the same general principles
that determine the form of the effective action.
In particular, the effective fields representing fields at positive
flow time have no O($a$) terms.

Such terms however occur in the case of the local
fields at flow time zero.
Considering again the theory with only massless quarks,
the effective fields representing the 
axial currents with flavour indices $r\neq s$, 
for example, are given by
\equation{
  (A_{\rm eff})_{\mu}^{rs}(x)
  =\psibar_r(x)\dirac{\mu}\dirac{5}\psi_s(x)+
   ac_5\partial_{\mu}\left\{\psibar_r(x)\dirac{5}\psi_s(x)\right\}
  \noenum
  \nexteq{2.0ex}
  {\phantom{(A_{\rm eff})_{\mu}^{rs}(x)=}}
  +ac_6\left\{\lambdabar_r(0,x)\dirac{\mu}\dirac{5}\psi_s(x)
  +\psibar_r(x)\dirac{\mu}\dirac{5}\lambda_s(0,x)
  \right\}+\ldots,
  \enum
}
where the ellipsis stands for terms of order $a^2$.
The associated effective pseudo-scalar densities,
\equation{
  (P_{\rm eff})^{rs}(x)=
  \psibar_r(x)\dirac{5}\psi_s(x)
  \noenum
  \noenum
  \nexteq{2.0ex}
  {\phantom{(P_{\rm eff})^{rs}(x)=}}
  +
  ac_7\left\{\lambdabar_r(0,x)\dirac{5}\psi_s(x)
  +\psibar_r(x)\dirac{5}\lambda_s(0,x)
  \right\}+\ldots,
  \enum
}
have a similar expansion.

The mass-dependent O($a$) corrections to these fields 
coincide with the ones in the theory in 4 dimensions
[\cite{ImpNonD}].
In general, the determination of the O($a$) 
terms however requires a careful consideration of a possible mixing
of fields and may lead to fairly complicated expressions.


\lussubsec 6.3\/ {\rm O($a$)\kern-0.5pt} improved lattice action

Following common practice, 
the mass-dependent O($a$) counterterms
are omitted in the improved action and are instead included
in the parameter and field renormalization factors [\cite{SFimp}].
The O($a$) counterterms
in 4+1 dimensions are then given by
\equation{
  \delta S_{\rm tot}=\sum_x\left\{\csw
  \psibar(x)\frac{i}{4}\sigma_{\mu\nu}\widehat{F}_{\mu\nu}(x)\psi(x)
  +\cfl\lambdabar(0,x)\lambda(0,x)\right\},
  \enum
}
where $\widehat{F}_{\mu\nu}(x)$ denotes the standard clover
expression for the gauge field tensor $F_{\mu\nu}(x)$ on the lattice.
While the first term in eq.~(6.7) (the Sheikholeslami--Wohlert term 
[\cite{SW},\cite{SFimp}]) is already required for the 
improvement of the theory in 4 dimensions, the one
proportional to the improvement
coefficient $\cfl(g_0)$
is needed to cancel the O($a$) lattice
effects in correlation functions involving the bulk fields.

Having specified the action,
the Wick contractions
of the basic fermion fields in the O($a$) improved theory
can be worked out straightforwardly 
(see appendix C).
In the continuous time limit, the contractions 
are practically the same
as in the continuum theory with the obvious modifications
(integrals over the space-time coordinates are replaced by
sums over lattice points, the quark propagator $S(x,y)$
by the inverse of the massive O($a$) improved lattice Dirac operator
and the kernel $K(t,x;s,y)$ by the fundamental solution of the 
quark flow equation on the lattice).
A special case is the contraction
\equation{
  \longwick{\chi(t,x)\chibar(s,y)}=\sum_{v,w}
  K(t,x;0,v)
  \left(S(v,w)-\cfl\delta_{vw}\right)
  K(s,y;0,w)^{\dagger}
  \enum 
}
of the time-dependent quark fields, which is the only one
depending on the improvement coefficient $\cfl$.
In particular, the contraction and thus the correlation functions
of the fields $\psi$ and $\psibar$ are independent of $\cfl$, as has to be
the case, since the theory in $4$ dimensions is already improved
through the inclusion of the Sheikholeslami--Wohlert term.

A perhaps puzzling aspect of eq.~(6.8) is the fact that the limit
\equation{
  \lim_{t\to0}\kern0.5pt\longwick{\chi(t,x)\chibar(s,y)}=
  \wick{\psi(x)\chibar(s,y)}-\cfl K(s,y;0,x)^{\dagger}
  \enum
}
differs from $\wick{\psi(x)\chibar(s,y)}$ by a 
term of order $a$, which is not just a contact term.
The subtractions required for
O($a$) improvement thus depend on whether the quark fields
in the correlation functions are at zero and 
non-zero flow time. This is actually not too surprising
since these fields also renormalize differently.
Both differences merely reflect the fact that 
the continuum limit of correlation
functions involving the time-dependent fields 
must be taken at fixed flow times given in physical units
and that the Symanzik effective theory in 4+1 dimensions
correctly describes the deviation of the lattice theory
from its continuum limit only when the latter is approached 
in this way.


\lussubsec 6.4 Renormalized improved fields

As usual, the additively renormalized bare quark masses 
$\mq{r}$ are defined by
\equation{
  \mq{r}=m_{0,r}-m_{\rm c},
  \enum
}
where $m_{\rm c}$ denotes the critical bare mass in the theory 
with mass-degenerate quarks.
It is also helpful to introduce
the associated subtracted bare mass matrix $\Mq$ and the combinations
\equation{
   \mq{rs}=\frac{1}{2}\left(\mq{r}+\mq{s}\right)
   \enum
}
of quark masses.

The lattice fermion fields at positive flow time require multiplicative 
renormalization and O($a$) improvement by a mass-dependent factor
(cf.~subsect.~6.1).
Quark and anti\-quark fields 
are scaled with the same factor,
\equation{
  \rens{\chi}{r}=\{Z_{\chi}(1+b_{\chi}\mq{r}+
  \bar{b}_{\chi}\tr\Mq)\}^{1/2}\chi_r,
  \enum
}
while the Lagrange-multiplier fields $\lambda,\lambdabar$ 
are renormalized with the inverse of the same factor%
\kern1pt\sfn{$\dagger$}{%
Following refs.~[\cite{SFimp},\cite{ImpNonD}],
$b_X$ and $\bar{b}_X$ generically denote improvement
coefficients multiplying mass-dependent O($a$) 
counterterms.}.
Time-dependent composite fields like the 
pseudo-scalar density 
\equation{
  \rens{P}{t}^{rs}=Z_{\chi}(
  1+b_{\chi}\mq{rs}+\bar{b}_{\chi}\tr\Mq)
  P_t^{rs}
  \enum
}
then renormalize according to their quark content.

Starting from the bare fields
\equation{
  A_{\mu}^{rs}(x)=\psibar_r(x)\dirac{\mu}\dirac{5}\psi_s(x),
  \enum
  \nexteq{2.5ex}
  \tilde{A}^{rs}_{\mu}(x)=
  \lambdabar_r(0,x)\dirac{\mu}\dirac{5}\psi_s(x)
  +\psibar_r(x)\dirac{\mu}\dirac{5}\lambda_s(0,x),
  \enum
  \nexteq{2.5ex}
  P^{rs}(x)=\psibar_r(x)\dirac{5}\psi_s(x),
  \enum
  \nexteq{2.5ex}
  \tilde{P}^{rs}(x)=
  \lambdabar_r(0,x)\dirac{5}\psi_s(x)
  +\psibar_r(x)\dirac{5}\lambda_s(0,x),
  \enum
}
the improved flavour non-singlet axial current and density are
given by
\equation{
  \imps{A}{\mu}^{rs}=A^{rs}_{\mu}
  +\ca\ring{\partial}_{\mu}P^{rs}
  +\cta\tilde{A}^{rs}_{\mu},
  \enum
  \nexteq{2.5ex}
  \imp{P}^{rs}=P^{rs}+\ctp\tilde{P}^{rs},
  \enum
}
and these are renormalized according to [\cite{SFimp},\cite{ImpNonD}]
\equation{
  \rens{A}{\mu}^{rs}=\ZA(1+b_A\mq{rs}+\bar{b}_A\tr\Mq)
  \imps{A}{\mu}^{rs},
  \enum
  \nexteq{2.5ex}
  \ren{P}^{rs}=\ZP(1+b_P\mq{rs}+\bar{b}_P\tr\Mq)
  \imp{P}^{rs}.
  \enum
}
The new improvement coefficients, $\cta$ and $\ctp$, are required for
the O($a$) improvement of correlation functions involving the fermion
fields at non-zero flow time, but all other coefficients already
occur in the theory in four dimensions [\cite{SW}--\cite{ImpNonD}].
At tree-level of perturbation theory,
\equation{
  \cfl=\frac{1}{2},
  \qquad
  \cta=\ctp=-\frac{1}{2},
  \qquad
  b_{\chi}=1.
  \enum
}
All coefficients $\bar{b}_X$ are, incidentally, of order $g_0^4$
[\cite{ImpNonD}].


\lussubsec 6.5 Improvement and renormalization of the field $\Ptax^{rs}$

In the chiral symmetry relations derived in 
sect.~4, the field $\Ptax^{rs}$ plays a prominent r\^ole.
The improvement and renormalization is a bit complicated 
in this case and is therefore discussed separately from the
other fields.

As already noted in subsect.~4.2, the field
can only mix with 
local composite fields that include at least one Lagrange-multiplier 
field. Charge conjugation, the lattice symmetries and the flavour symmetry
then imply that the field is multiplicatively renormalizable. 
There are, however, several fields that can mix 
with the density at order $a$,
among them two fields
\equation{
  \hat{P}^{rs}(x)=\lambdabar_r(0,x)\dirac{5}\lambda_s(0,x),
  \enum
  \nexteq{2.5ex}
  Q^{rs}(x)=\lambdabar_r(0,x)\dirac{5}\psi_s(x)-
  \psibar_r(x)\dirac{5}\lambda_s(0,x),
  \enum
}
that have not appeared before.
Inspection then shows that 
a possible choice of the improved and renormalized densities is
\equation{
  \imp{\Ptax}^{rs}=
  \Ptax^{rs}+\hat{c}_A
  \ring{\partial}_{\mu}
  \tilde{A}_{\mu}^{rs}+\hat{c}_P\hat{P}^{rs},
  \enum
  \nexteq{2.5ex}
  \ren{\Ptax}^{rs}=
  \tilde{Z}_P\bigl\{
  (1+b_{\Ptax}\mq{rs}+\bar{b}_{\Ptax}\tr\Mq)
  \imp{\Ptax}^{rs}
  +\hat{b}_{\Ptax}(\mq{r}-\mq{s})Q^{rs}\bigr\},
  \enum
}
where, as usual, the field equations were used to reduce the
number of terms.
The fields on the right of these equations 
are distinguished by their symmetry properties or content in
Lagrange-multiplier fields.
They are all multiplicatively renormalizable and the contributions
of the counterterms proportional to $\tilde{A}_{\mu}^{rs}$, $\hat{P}^{rs}$ and
$Q^{rs}$ in on-shell correlation functions are therefore of order $a$.

The expression (6.25),(6.26) for the renormalized improved density
can be simplified by noting that the relations
\equation{
  \sum_x\langle\Ptax^{rs}(x)P^{sr}_t(y)\rangle
  =\langle S_t^{rr}(y)\rangle+\langle S_t^{ss}(y)\rangle
  +2\cfl\sum_x\langle\hat{P}^{rs}(x)P^{sr}_t(y)\rangle,
  \enum
  \nexteq{2.5ex}
  \sum_x\langle Q^{rs}(x)P^{sr}_t(y)\rangle
  =\langle S_t^{ss}(y)\rangle-\langle S_t^{rr}(y)\rangle,
  \enum
}
hold exactly, for any $t>0$, as a consequence of
the form of the Wick contractions of the fermion fields. 
On the other hand, as discussed
in subsect.~4.2, the renormalization
constants can be (and are to be) chosen so that the normalization 
condition (4.15) is satisfied in the continuum limit.
Since the correlation functions in eq.~(4.15) converge to the
continuum limit with a rate proportional to $a^2$,
the comparison with the unrenormalized identities, eqs.~(6.27),(6.28),
then shows that
\equation{
  \tilde{Z}_P=1,
  \enum
  \nexteq{2.5ex}
  \hat{c}_P=-2\cfl,
  \quad
  b_{\Ptax}=\bar{b}_{\Ptax}=0,
  \quad
  \hat{b}_{\Ptax}=-\frac{1}{2}b_{\chi}.
  \enum
}
The renormalized density 
\equation{
  \ren{\Ptax}^{rs}=
  \Ptax^{rs}+
  \hat{c}_A
  \ring{\partial}_{\mu}
  \tilde{A}_{\mu}^{rs}-2\cfl\hat{P}^{rs}
  -b_{\chi}\frac{1}{2}(\mq{r}-\mq{s})Q^{rs}
  \enum
}
thus assumes a fairly simple form,
in which $\hat{c}_A$ is the only new improvement coefficient.


\lussubsec 6.6 PCAC relation on the lattice

In the continuum limit, the renormalized PCAC relation (4.14) holds 
provided the renormalization constants $Z_A,Z_P$ and the renormalized
quark masses $\mR{r}$ are chosen appropriately.
If also all improvement coefficients are properly tuned, this
implies
\equation{
  \langle\partial\ren{A}^{rs}(x)\rens{P}{t}^{sr}(y)\rangle
  \noenum
  \nexteq{2.0ex}
  \qquad
  =(\mR{r}+\mR{s})\langle\ren{P}^{rs}(x)\rens{P}{t}^{sr}(y)\rangle
  -\langle\ren{\Ptax}^{rs}(x)\rens{P}{t}^{sr}(y)\rangle
  +\rmO(a^2),
  \enum
}
where, following the tradition, the divergence of the 
improved axial current is taken to be
\equation{
  \partial\ren{A}^{rs}=\ZA\left(1+b_A\mq{rs}+\bar{b}_A\tr\Mq\right)
  \partial\imp{A}^{rs},
  \enum
  \nexteq{2.5ex}
  \partial\imp{A}^{rs}=
  \ring{\partial}_{\mu}\{A^{rs}_{\mu}+
  \cta\tilde{A}^{rs}_{\mu}\}
  +\ca\drvstar{\mu}\drv{\mu}P^{rs}.
  \enum
}
A technical detail worth emphasizing here again is the fact that
the PCAC relation (6.32) holds at all $x$, including $x=y$, as long as 
the flow time $t$ is set to a positive value given in physical units.
